\documentclass[12pt, a4paper]{article}

% ─── Packages ─────────────────────────────────────────────────
\usepackage{fontspec}
\usepackage{polyglossia}
\setmainlanguage{vietnamese}
\setmainfont{DejaVu Serif}[
  BoldFont       = DejaVu Serif Bold,
  ItalicFont     = DejaVu Serif Italic,
  BoldItalicFont = DejaVu Serif Bold Italic
]

\usepackage[margin=2.5cm]{geometry}
\usepackage{amsmath, amssymb}
\usepackage{graphicx}
\usepackage{booktabs}
\usepackage{multirow}
\usepackage{array}
\usepackage{xcolor}
\usepackage[unicode, pdfencoding=auto]{hyperref}
\usepackage{caption}
\usepackage{subcaption}
\usepackage{setspace}
\usepackage{fancyhdr}
\usepackage{titlesec}
\usepackage{enumitem}
\usepackage{float}
\usepackage{makecell}
\usepackage[numbers, sort&compress]{natbib}
\usepackage{microtype}

% ─── Hyperref setup ────────────────────────────────────────────
\hypersetup{
  colorlinks = true,
  linkcolor  = blue,
  citecolor  = blue,
  urlcolor   = blue,
  pdftitle   = {Thuc nghiem CVRP + xAI},
  pdfauthor  = {},
}

% ─── Header / page style ───────────────────────────────────────
\setlength{\headheight}{14pt}
\pagestyle{fancy}
\fancyhf{}
\rhead{\small Thực nghiệm CVRP + xAI}
\rfoot{\thepage}

% ─── Spacing & graphics ────────────────────────────────────────
\onehalfspacing
\graphicspath{{figures/}}

% ─── Overfull tolerance ────────────────────────────────────────
\setlength{\emergencystretch}{3em}
\tolerance 9999
\hbadness 9999

% ─── Màu sắc model ────────────────────────────────────────────
\definecolor{colorDL}{HTML}{2196F3}
\definecolor{colorSF}{HTML}{4CAF50}
\definecolor{colorKN}{HTML}{FF5722}
\definecolor{colorM2F}{HTML}{9C27B0}

% ─── Caption style ─────────────────────────────────────────────
\captionsetup{font=small, labelfont=bf}
\captionsetup[subfigure]{font=small}

% ═══════════════════════════════════════════════════════════════
\begin{document}

% ─── Trang bìa ────────────────────────────────────────────────
\begin{titlepage}
\centering
\vspace*{2cm}
{\LARGE \textbf{Thực nghiệm tái tạo và ứng dụng\\[0.3em]
Trí tuệ nhân tạo giải thích được (xAI)\\[0.3em]
trên bộ dữ liệu phân đoạn bông lúa CVRP}}\\[1.5cm]
{\large Báo cáo thực nghiệm}\\[1cm]

\rule{0.6\linewidth}{0.5pt}\\[1cm]

\begin{tabular}{@{}ll}
\textbf{Bộ dữ liệu gốc:}  & CVRP (Tang et al., \textit{Plant Phenomics} 2025)\\[4pt]
\textbf{Phương pháp:}      & Semantic Segmentation + Explainable AI\\[4pt]
\textbf{Môi trường:}       & PyTorch 2.1 · MMSeg 1.2.2 · NVIDIA RTX 2080 Ti\\[4pt]
\textbf{Năm:}              & 2026\\
\end{tabular}

\vfill
\end{titlepage}

% ─── Mục lục ──────────────────────────────────────────────────
\tableofcontents
\listoffigures
\listoftables
\newpage

% ═══════════════════════════════════════════════════════════════
\section{Giới thiệu}
\label{sec:intro}

Phân đoạn bông lúa (\textit{rice panicle segmentation}) là bài toán quan trọng trong kiểu hình học thực vật (\textit{plant phenomics}), phục vụ ước lượng năng suất và nghiên cứu chọn giống. Bài báo gốc của Tang et al.~\cite{tang2025cvrp} giới thiệu bộ dữ liệu CVRP --- bộ dữ liệu đa giống, đa góc nhìn gồm 2.426 ảnh cây lúa với chú thích chất lượng cao, và đánh giá bốn mô hình phân đoạn ngữ nghĩa tiên tiến.

Báo cáo này thực hiện hai đóng góp chính:
\begin{enumerate}[label=(\arabic*)]
  \item \textbf{Tái tạo thực nghiệm}: Huấn luyện và đánh giá lại bốn mô hình gồm DeepLabV3+, SegFormer, K-Net và Mask2Former trên bộ dữ liệu CVRP, so sánh định lượng với kết quả trong bài báo gốc và phân tích nguyên nhân sai lệch.
  \item \textbf{Ứng dụng xAI}: Áp dụng các phương pháp Trí tuệ nhân tạo giải thích được (cả \textit{saliency-based} lẫn \textit{counterfactual}) để phân tích cơ chế hoạt động của từng mô hình trên bài toán phân đoạn bông lúa.
\end{enumerate}

% ═══════════════════════════════════════════════════════════════
\section{Bộ dữ liệu CVRP}
\label{sec:dataset}

\subsection{Mô tả tổng quan}

Bộ dữ liệu CVRP~\cite{tang2025cvrp} được thu thập từ ruộng lúa nước tại Nam Kinh, Trung Quốc với 231 giống địa phương và 50 giống hiện đại. Cấu trúc dữ liệu gồm ba phần như Bảng~\ref{tab:dataset}.

\begin{table}[H]
\centering
\caption{Thống kê bộ dữ liệu CVRP}
\label{tab:dataset}
\begin{tabular}{lcccc}
\toprule
\textbf{Loại dữ liệu} & \makecell{\textbf{Số}\\\textbf{giống}} & \makecell{\textbf{Số}\\\textbf{ảnh}} & \makecell{\textbf{Anno-}\\\textbf{tation}} & \textbf{Kích thước} \\
\midrule
Multi-cultivar & 282 & 2.205 & \checkmark & $1080\times1920$ \\
Field scenes   & --- & 100   & \checkmark & Biến đổi \\
Indoor (lab)   & 41  & 123   & $\times$   & $1500\times2007$ \\
\midrule
\textbf{Tổng} & 282 & 2.428 & 2.305 & --- \\
\bottomrule
\end{tabular}
\end{table}

Annotation được lưu dạng Palette PNG hai giá trị: \texttt{0}~=~nền, \texttt{1}~=~bông lúa. Tỉ lệ pixel bông chỉ \textbf{7\%} trên tập train/val (Hình~\ref{fig:dataset_stats}b), cho thấy mất cân bằng lớp đáng kể.

\subsection{Phân chia dữ liệu}

\begin{itemize}
  \item \textbf{Train}: 2.116 ảnh (multi-cultivar, 282 giống)
  \item \textbf{Val}: 80 ảnh (multi-cultivar) --- dùng để tái tạo Table~2 của paper
  \item \textbf{Field test}: 100 ảnh field scenes --- dùng để tái tạo Table~3 của paper
\end{itemize}

\subsection{Phân tích bộ dữ liệu}

\begin{figure}[H]
  \centering
  \begin{subfigure}[b]{0.48\textwidth}
    \includegraphics[width=\textwidth]{01_dataset_analysis/fig_multiview_sample.png}
    \caption{Bốn góc nhìn của cùng một giống}
    \label{fig:multiview}
  \end{subfigure}
  \hfill
  \begin{subfigure}[b]{0.48\textwidth}
    \includegraphics[width=\textwidth]{01_dataset_analysis/fig_annotation_examples.png}
    \caption{Ví dụ annotation (mask đỏ = bông)}
    \label{fig:ann_examples}
  \end{subfigure}
  \caption{Đặc điểm bộ dữ liệu CVRP: đa góc nhìn và annotation chi tiết.}
  \label{fig:dataset_examples}
\end{figure}

\begin{figure}[H]
  \centering
  \begin{subfigure}[b]{0.48\textwidth}
    \includegraphics[width=\textwidth]{01_dataset_analysis/fig_images_per_cultivar.png}
    \caption{Phân bố số ảnh theo giống}
    \label{fig:imgs_per_cultivar}
  \end{subfigure}
  \hfill
  \begin{subfigure}[b]{0.48\textwidth}
    \includegraphics[width=\textwidth]{01_dataset_analysis/fig_annotation_ratio_boxplot.png}
    \caption{Tỉ lệ pixel bông theo tập dữ liệu}
    \label{fig:ann_ratio}
  \end{subfigure}
  \caption{Thống kê phân bố: trung bình 7--15\% pixel bông, mất cân bằng lớp đáng kể.}
  \label{fig:dataset_stats}
\end{figure}

% ═══════════════════════════════════════════════════════════════
\section{Phương pháp phân đoạn ngữ nghĩa}
\label{sec:method}

\subsection{Bốn mô hình được đánh giá}

Chúng tôi sử dụng MMSegmentation v1.2.2~\cite{mmseg2020} để triển khai và huấn luyện bốn mô hình tiên tiến:

\subsubsection{DeepLabV3+ (ResNet-101)}
DeepLabV3+~\cite{chen2018deeplabv3plus} dùng Atrous Spatial Pyramid Pooling (ASPP) để nắm bắt ngữ cảnh đa tỷ lệ, kết hợp với encoder-decoder để đạt độ phân giải dự đoán cao.

\subsubsection{SegFormer (MiT-B2)}
SegFormer~\cite{xie2021segformer} kết hợp Mix Vision Transformer (MiT) phân cấp với bộ giải mã MLP nhẹ, không cần positional encoding, thích ứng với kích thước ảnh tùy ý.

\subsubsection{K-Net (Swin-B)}
K-Net~\cite{zhang2021knet} dùng kernel động cập nhật iterative, mỗi kernel học cách attend vào vùng đối tượng khác nhau. Chúng tôi dùng Swin-B làm backbone (trong khi paper gốc không nêu rõ backbone K-Net).

\subsubsection{Mask2Former (Swin-B)}
Mask2Former~\cite{cheng2022mask2former} dựa trên masked attention transformer decoder với pixel decoder MSDeformAttn, đạt hiệu suất tốt nhất trong cả bốn mô hình.

\subsection{Cấu hình huấn luyện}

\begin{table}[H]
\centering
\caption{Hyperparameter huấn luyện --- đồng nhất theo cấu hình bài báo gốc}
\label{tab:training_config}
\begin{tabular}{@{}ll@{}}
\toprule
\textbf{Thông số} & \textbf{Giá trị} \\
\midrule
Optimizer         & AdamW \\
Learning rate     & $10^{-4}$ \\
Backbone LR       & $10^{-5}$ (lr\_mult = 0.1) \\
Batch size        & 2 \\
Số iterations     & 80.000 \\
Crop size         & $512 \times 512$ \\
LR warmup         & Linear, 1.500 iter đầu \\
GPU               & NVIDIA RTX 2080 Ti (11\,GB VRAM) \\
Class weight      & [0.3, 3.0] (background, panicle) \\
\bottomrule
\end{tabular}
\end{table}

Backbone pretrained: ResNet-101 (ImageNet-1K), MiT-B2 (ImageNet-1K), Swin-B (ImageNet-22K).

% ═══════════════════════════════════════════════════════════════
\section{Kết quả phân đoạn}
\label{sec:results}

\subsection{Kết quả định lượng}

\begin{table}[H]
\centering
\caption{Panicle IoU trên tập Val --- tái tạo Table~2 của \cite{tang2025cvrp}}
\label{tab:results_val}
\begin{tabular}{lcccccc}
\toprule
\multirow{2}{*}{\textbf{Model}} &
\multicolumn{3}{c}{\textbf{Kết quả của chúng tôi}} &
\multicolumn{2}{c}{\textbf{Paper gốc}} &
\multirow{2}{*}{$\Delta$IoU} \\
\cmidrule(lr){2-4}\cmidrule(lr){5-6}
 & IoU & Recall & mIoU & IoU & Recall & \\
\midrule
{\color{colorDL}DeepLabV3+}  & 65.66 & 98.16 & 80.93 & 71.02 & 86.55 & $-5.36$ \\
{\color{colorSF}SegFormer}    & 66.61 & 97.66 & 81.49 & 71.19 & 88.11 & $-4.58$ \\
{\color{colorKN}K-Net}        & 75.56 & 88.44 & 86.73 & 72.53 & 85.85 & $\mathbf{+3.03}$ \\
{\color{colorM2F}Mask2Former} & \textbf{77.58} & 89.36 & \textbf{87.84} & \textbf{73.38} & 88.49 & $\mathbf{+4.20}$ \\
\midrule
MAE & \multicolumn{6}{c}{4.29\%} \\
\bottomrule
\end{tabular}
\end{table}

\begin{table}[H]
\centering
\caption{Panicle IoU trên tập Field Test --- tái tạo Table~3 của \cite{tang2025cvrp}}
\label{tab:results_field}
\begin{tabular}{lcccccc}
\toprule
\multirow{2}{*}{\textbf{Model}} &
\multicolumn{3}{c}{\textbf{Kết quả của chúng tôi}} &
\multicolumn{2}{c}{\textbf{Paper gốc}} &
\multirow{2}{*}{$\Delta$IoU} \\
\cmidrule(lr){2-4}\cmidrule(lr){5-6}
 & IoU & Recall & mIoU & IoU & Recall & \\
\midrule
{\color{colorDL}DeepLabV3+}  & 63.46 & 93.94 & 77.03 & 62.82 & 74.18 & $+0.64$ \\
{\color{colorSF}SegFormer}    & 63.12 & 95.21 & 76.72 & 68.24 & 84.08 & $-5.12$ \\
{\color{colorKN}K-Net}        & 69.23 & 78.81 & 81.65 & 67.34 & 77.50 & $+1.89$ \\
{\color{colorM2F}Mask2Former} & \textbf{72.00} & 82.53 & \textbf{83.27} & \textbf{76.56} & 84.60 & $-4.56$ \\
\midrule
MAE & \multicolumn{6}{c}{3.05\%} \\
\bottomrule
\end{tabular}
\end{table}

Sai số tuyệt đối trung bình (MAE) đạt \textbf{4.29\%} (val) và \textbf{3.05\%} (field test) --- nằm trong ngưỡng chấp nhận ($\leq 5\%$) cho thực nghiệm tái tạo. Thứ hạng ranking \textbf{Mask2Former > K-Net > SegFormer > DeepLabV3+} hoàn toàn nhất quán với paper gốc.

\subsection{Phân tích sai lệch}

\textbf{(1) Class weight gây mất cân bằng Precision--Recall.}
Chúng tôi dùng \texttt{class\_weight=[0.3,\,3.0]} để xử lý imbalanced data (bông chỉ chiếm 7\% pixel), trong khi paper gốc không đề cập class weight. Hệ quả: Recall của chúng tôi đạt 93--98\% (vs.\ 74--88\% của paper), còn Precision chỉ 59--71\% (vs.\ 74--89\%). IoU thấp hơn ở CNN-based models do FP tăng cao.

\textbf{(2) K-Net và Mask2Former vượt paper} nhờ: (a)~nhiều training data hơn (2.116 vs.\ 727 ảnh); (b)~data leakage ở mức giống (val split random khiến tất cả 67 giống trong val đều có ảnh trong train); (c)~K-Net dùng Swin-B thay vì Swin-T.

\subsection{Đường cong huấn luyện}

\begin{figure}[H]
  \centering
  \includegraphics[width=\textwidth]{03_segmentation_curves/fig_training_curves.png}
  \caption{Loss huấn luyện và Val~mIoU theo iteration. Tất cả mô hình hội tụ ổn định sau 80k iter. Mask2Former có loss tổng cao (do tổng hợp loss cls + mask + dice qua 9 decoder stages) nhưng đạt mIoU cao nhất.}
  \label{fig:training_curves}
\end{figure}

\begin{figure}[H]
  \centering
  \begin{subfigure}[b]{0.53\textwidth}
    \includegraphics[width=\textwidth]{03_segmentation_curves/fig_results_comparison.png}
    \caption{So sánh Panicle IoU với paper}
  \end{subfigure}
  \hfill
  \begin{subfigure}[b]{0.44\textwidth}
    \includegraphics[width=\textwidth]{03_segmentation_curves/fig_radar_comparison.png}
    \caption{Radar chart đa metric}
  \end{subfigure}
  \caption{So sánh kết quả. (a)~K-Net và Mask2Former vượt paper trên val; Mask2Former yếu hơn trên field test do domain gap. (b)~Radar chart so sánh tổng thể cả hai split.}
  \label{fig:results_comparison}
\end{figure}

\subsection{Kết quả định tính}

\begin{figure}[H]
  \centering
  \includegraphics[width=\textwidth]{02_segmentation_qualitative/fig_4models_grid.png}
  \caption{So sánh trực quan bốn mô hình trên val set. Hàng đầu: ảnh gốc. Hàng 2--5: kết quả phân đoạn overlay (xanh lá = TP, đỏ = FP, vàng = FN). Mask2Former cho boundary rõ nhất và ít FP nhất.}
  \label{fig:4models_grid}
\end{figure}

\begin{figure}[H]
  \centering
  \begin{subfigure}[b]{0.48\textwidth}
    \includegraphics[width=\textwidth]{02_segmentation_qualitative/fig_success_cases.png}
    \caption{Trường hợp thành công (val set)}
  \end{subfigure}
  \hfill
  \begin{subfigure}[b]{0.48\textwidth}
    \includegraphics[width=\textwidth]{02_segmentation_qualitative/fig_failure_cases.png}
    \caption{Trường hợp thất bại (field test)}
  \end{subfigure}
  \caption{(a)~Mask2Former phân đoạn chính xác trên val set. (b)~Các trường hợp thách thức trên field test: cỏ dại bị nhầm thành bông, lá lúa che khuất bông, và bông nhỏ bị bỏ sót.}
  \label{fig:qualitative}
\end{figure}

% ═══════════════════════════════════════════════════════════════
\section{Phân đoạn hạt lúa với SAM}
\label{sec:sam}

Chúng tôi đánh giá Segment Anything Model~\cite{kirillov2023sam} (SAM ViT-H) trên 82 ảnh indoor (loại bỏ suffix \_3 không rõ trạng thái), tái tạo Section~3.3 của bài báo gốc.

\begin{table}[H]
\centering
\caption{Kết quả đếm hạt lúa bằng SAM ViT-H}
\label{tab:sam_results}
\begin{tabular}{lccccc}
\toprule
\textbf{Trạng thái} & \textbf{N ảnh} & \textbf{Avg hạt/ảnh} & \textbf{Std} & \textbf{Min} & \textbf{Max} \\
\midrule
Tự nhiên (suffix \_1) & 41 & 73.6 & 20.2 & 31 & 118 \\
Mở ra (suffix \_2)    & 41 & 67.3 & 20.1 & 12 & 110 \\
\bottomrule
\end{tabular}
\end{table}

\noindent\textbf{Lưu ý:} Paper gốc báo cáo tỉ lệ phần trăm hạt phát hiện (73\% natural, 83\% unfolded) so với đếm tay. Chúng tôi không có ground truth đếm tay nên chỉ báo số lượng tuyệt đối SAM phát hiện được.

\begin{figure}[H]
  \centering
  \begin{subfigure}[b]{0.48\textwidth}
    \includegraphics[width=\textwidth]{04_sam_grains/fig_sam_natural.png}
    \caption{Bông tự nhiên (suffix \_1)}
  \end{subfigure}
  \hfill
  \begin{subfigure}[b]{0.48\textwidth}
    \includegraphics[width=\textwidth]{04_sam_grains/fig_sam_unfolded.png}
    \caption{Bông mở ra (suffix \_2)}
  \end{subfigure}
  \caption{SAM ViT-H phân đoạn hạt lúa trên ảnh indoor nền đen. Mỗi màu đại diện một hạt được phát hiện.}
  \label{fig:sam_grains}
\end{figure}

\begin{figure}[H]
  \centering
  \includegraphics[width=0.65\textwidth]{04_sam_grains/fig_sam_comparison.png}
  \caption{So sánh số hạt phát hiện theo trạng thái bông và theo giống lúa. Độ biến thiên cao (std $\approx$~20) phản ánh sự đa dạng về đặc điểm bông giữa các giống.}
  \label{fig:sam_comparison}
\end{figure}

% ═══════════════════════════════════════════════════════════════
\section{Tái tạo 3D cây lúa}
\label{sec:3d}

\begin{figure}[H]
  \centering
  \includegraphics[width=0.8\textwidth]{05_reconstruction/fig_3d_reconstruction_grid.png}
  \caption{Tái tạo 3D ba giống lúa (G1204, G1210, G1224) từ file \texttt{.ply} có sẵn trong dataset. Cột trái: point cloud gốc (màu sắc thực). Cột phải: bông lúa được highlight màu vàng bằng phân tích màu sắc. Render bằng Open3D và matplotlib.}
  \label{fig:3d_recon}
\end{figure}

% ═══════════════════════════════════════════════════════════════
\section{Ứng dụng Explainable AI (xAI)}
\label{sec:xai}

Đây là phần đóng góp mới so với bài báo gốc. Chúng tôi áp dụng xAI để trả lời năm câu hỏi nghiên cứu:

\begin{description}[leftmargin=3.5cm, labelwidth=3cm]
  \item[\textbf{RQ1}] Model focus vào vùng nào khi predict bông lúa?
  \item[\textbf{RQ2}] CNN và Transformer có attention pattern khác nhau không?
  \item[\textbf{RQ3}] Tại sao model thất bại trên field scenes?
  \item[\textbf{RQ4}] Chất lượng explanation cao đến mức nào? (định lượng)
  \item[\textbf{RQ5}] Thay đổi tối thiểu nào khiến model đổi dự đoán? (counterfactual)
\end{description}

\subsection{Phương pháp xAI --- Saliency-based}

\subsubsection{GradCAM++ \cite{chattopadhay2018gradcampp}}
GradCAM++ cải tiến GradCAM bằng gradient bậc hai, highlight tốt hơn khi có nhiều đối tượng cùng loại. Áp dụng cho \textit{DeepLabV3+}, \textit{SegFormer} và \textit{K-Net}.

\smallskip
\noindent\textbf{Target layers:}
\begin{itemize}[noitemsep]
  \item DeepLabV3+: \texttt{backbone.layer4[-1]} (Bottleneck cuối ResNet-101)
  \item SegFormer: \texttt{decode\_head.fusion\_conv} (ConvModule tổng hợp 4-scale features)
  \item K-Net: \texttt{decode\_head.kernel\_generate\_head.fpn\_bottleneck}
\end{itemize}

\subsubsection{EigenCAM \cite{muhammad2020eigencam}}
EigenCAM không dùng gradient --- áp dụng SVD trên activation tensor, lấy principal component đầu tiên làm heatmap. Phù hợp với Transformer vì gradient của ViT/Swin thường noisy. Áp dụng cho \textit{Mask2Former} với reshape transform $(B, L, C) \to (B, C, H, W)$.

\subsubsection{Energy Pointing Game (EPG) \cite{zhang2018top}}
Metric định lượng chất lượng explanation:
\begin{equation}
\text{EPG} = \frac{\displaystyle\sum_{i \in \text{GT\_mask}} \text{CAM}(i)}{\displaystyle\sum_{i} \text{CAM}(i)}
\end{equation}
EPG cao $\Rightarrow$ CAM energy tập trung đúng vùng bông lúa (GT mask).

\begin{figure}[H]
  \centering
  \includegraphics[width=\textwidth]{06_xai/gradcam/fig_cam_4models_grid.png}
  \caption{\textbf{RQ1 + RQ2}: GradCAM++ (DeepLabV3+, SegFormer, K-Net) và EigenCAM (Mask2Former) trên 4 ảnh val set. Vùng màu nóng (đỏ/vàng) cho thấy vùng model chú ý khi dự đoán bông. EPG score hiển thị trên mỗi ô.}
  \label{fig:cam_4models}
\end{figure}

\begin{figure}[H]
  \centering
  \includegraphics[width=\textwidth]{06_xai/comparison/fig_cnn_vs_transformer.png}
  \caption{\textbf{RQ2}: So sánh trực tiếp GradCAM++ (DeepLabV3+/CNN) và EigenCAM (Mask2Former/Transformer). CNN tạo heatmap cục bộ, sắc nét tập trung vào texture bông; Transformer tạo heatmap lan rộng nắm bắt ngữ cảnh toàn ảnh. Đường đỏ là contour CAM tại ngưỡng~0.4.}
  \label{fig:cnn_vs_transformer}
\end{figure}

\begin{figure}[H]
  \centering
  \includegraphics[width=0.65\textwidth]{06_xai/quantitative/fig_cam_quantitative.png}
  \caption{\textbf{RQ4}: Energy Pointing Game score cho bốn mô hình. DeepLabV3+ đạt cao nhất (Val~EPG~=~0.167, Field~EPG~=~0.234), xác nhận CNN có attention cục bộ tập trung hơn Transformer vào vùng bông.}
  \label{fig:epg_quantitative}
\end{figure}

\begin{table}[H]
\centering
\caption{Kết quả định lượng xAI --- Energy Pointing Game Score (mean $\pm$ std)}
\label{tab:epg_results}
\begin{tabular}{@{}llcccc@{}}
\toprule
\textbf{Model} & \textbf{Method} &
\textbf{Val EPG} & \textbf{Std} &
\textbf{Field EPG} & \textbf{Std} \\
\midrule
{\color{colorDL}DeepLabV3+}  & GradCAM++ & \textbf{0.167} & 0.079 & \textbf{0.234} & 0.018 \\
{\color{colorSF}SegFormer}    & GradCAM++ & 0.094 & 0.062 & 0.157 & 0.049 \\
{\color{colorKN}K-Net}        & GradCAM++ & 0.062 & 0.044 & 0.100 & 0.034 \\
{\color{colorM2F}Mask2Former} & EigenCAM  & 0.073 & 0.070 & 0.147 & 0.121 \\
\bottomrule
\end{tabular}
\end{table}

\begin{figure}[H]
  \centering
  \includegraphics[width=\textwidth]{06_xai/failure/fig_failure_cam.png}
  \caption{\textbf{RQ3}: CAM trên failure cases của Mask2Former (field test). Hàng 1: ảnh gốc. Hàng 2: Ground Truth. Hàng 3: CAM heatmap. Hàng 4: GT (xanh lá) và FP region (đỏ). Model attend vào cỏ dại màu vàng-nâu tương tự bông (case 1), lá lúa giai đoạn chín (case 3) và các vùng xa bông (case 4).}
  \label{fig:failure_cam}
\end{figure}

\subsection{Phương pháp xAI --- Counterfactual}

\subsubsection{Occlusion Counterfactual Map \cite{zeiler2014visualizing}}
Slide patch $32\times32$ (stride~16) qua ảnh, thay bằng màu mean của ImageNet. Đo độ sụt giảm dự đoán:
\begin{equation}
\text{CF}(r,c) = \max\!\bigl(0,\;S_{\text{base}} - S_{\text{occ}(r,c)}\bigr)
\end{equation}
Vùng CF cao $\Rightarrow$ \textit{counterfactually necessary}: xóa đi thì model không còn detect bông.

\subsubsection{Deletion/Insertion AUC \cite{samek2017evaluating}}
Sắp xếp pixel theo độ quan trọng (từ CAM map), sau đó:
\begin{itemize}[noitemsep]
  \item \textbf{Deletion}: xóa dần pixel quan trọng $\to$ score giảm nhanh = CAM tốt
  \item \textbf{Insertion}: thêm dần pixel vào ảnh blank $\to$ score tăng nhanh = CAM tốt
\end{itemize}
AUC của hai curve đo chất lượng explanation.

\subsubsection{Superpixel Counterfactual \cite{petsiuk2018rise}}
Phân ảnh thành 40 superpixels (SLIC~\cite{petsiuk2018rise}). Với mỗi superpixel: xóa và đo drop IoU. Superpixel có drop lớn $\Rightarrow$ \textit{critical} cho detection.

\begin{figure}[H]
  \centering
  \includegraphics[width=\textwidth]{06_xai/counterfactual/fig_occlusion_cf_map.png}
  \caption{\textbf{RQ5 --- Phương pháp 1}: Occlusion Counterfactual Map cho ba mô hình trên bốn ảnh (2 val + 2 field). Vùng màu nóng = xóa vùng này thì model mất khả năng detect bông. CF-Align là tỉ lệ vùng counterfactual trùng với GT mask.}
  \label{fig:occlusion_cf}
\end{figure}

\begin{figure}[H]
  \centering
  \includegraphics[width=\textwidth]{06_xai/counterfactual/fig_deletion_insertion.png}
  \caption{\textbf{RQ5 --- Phương pháp 2}: Deletion/Insertion curves cho DeepLabV3+ trên hai ảnh val mẫu. Đường đỏ = deletion (xóa pixel quan trọng $\to$ score giảm). Đường xanh = insertion (thêm pixel $\to$ score tăng). Del~AUC~=~45.37; Ins~AUC~=~43.59. \textit{Ghi chú: Phương pháp này chỉ hoàn thành với DeepLabV3+ do giới hạn thời gian tính toán.}}
  \label{fig:del_ins}
\end{figure}

\begin{figure}[H]
  \centering
  \includegraphics[width=\textwidth]{06_xai/counterfactual/fig_superpixel_cf.png}
  \caption{\textbf{RQ5 --- Phương pháp 3}: Superpixel Counterfactual (Mask2Former). Hàng 1: ảnh gốc + GT. Hàng 2: superpixels SLIC. Hàng 3: critical superpixels được highlight (xanh lá = critical và đúng vùng bông; đỏ = critical nhưng nằm ngoài GT = model nhầm vùng này là quan trọng).}
  \label{fig:superpixel_cf}
\end{figure}

\subsection{Phân tích và thảo luận}

\subsubsection{CNN vs.\ Transformer --- Cơ chế attention khác nhau (RQ2)}
DeepLabV3+ (CNN) có EPG cao nhất (0.167) và CF-Align tốt nhất, xác nhận attention \textit{cục bộ} tập trung đúng texture bông lúa. Transformer (Mask2Former, K-Net) có EPG thấp hơn vì attention \textit{rộng} nắm bắt context toàn ảnh --- phù hợp để phân biệt bông/nền phức tạp nhưng khó giải thích theo vùng.

\subsubsection{Failure case analysis (RQ3)}
Từ Hình~\ref{fig:failure_cam}:
\begin{itemize}[noitemsep]
  \item \textbf{Cỏ dại}: Model attend vào texture vàng-nâu giống bông $\to$ FP tăng cao
  \item \textbf{Domain gap}: Plastic boards giả tạo trong training không có trong field test $\to$ generalization kém
  \item \textbf{Bông nhỏ}: CAM energy phân tán, không đủ tập trung để detect bông nhỏ $\to$ FN tăng
\end{itemize}

\subsubsection{Giới hạn kỹ thuật của xAI với Mask2Former}
Kiến trúc query-based khiến \texttt{decode\_head.forward()} yêu cầu \texttt{batch\_data\_samples} --- không tương thích với gradient hooks thông thường. Giải pháp tạm thời: dùng backbone wrapper (Swin-B $\to$ linear projection 2 classes). Đây là hạn chế kỹ thuật cần nghiên cứu tiếp với phương pháp như DINO attention~\cite{liu2021swin} hay Transformer rollout.

% ═══════════════════════════════════════════════════════════════
\section{Kết luận}
\label{sec:conclusion}

\subsection{Tổng kết đóng góp}

\begin{enumerate}
  \item \textbf{Tái tạo thành công}: MAE~=~4.29\% (val) và 3.05\% (field) --- trong ngưỡng chấp nhận được. Thứ hạng ranking 100\% nhất quán với paper gốc.
  \item \textbf{K-Net và Mask2Former vượt paper} (+3--4\% IoU) nhờ training data gấp 2.9 lần và backbone Swin-B.
  \item \textbf{xAI làm rõ cơ chế}: CNN attention cục bộ và tập trung; Transformer attention rộng nắm bắt ngữ cảnh.
  \item \textbf{Counterfactual xAI xác định vùng critical}: Occlusion Map và Superpixel CF chỉ ra vùng nào là \textit{counterfactually necessary} để model detect bông.
  \item \textbf{Failure analysis}: Domain gap (plastic boards), texture ambiguity (cỏ dại/bông) và bông nhỏ là ba nguyên nhân chính gây sai sót.
\end{enumerate}

\subsection{Hướng phát triển}

\begin{itemize}
  \item Bỏ hoặc giảm class weight về [0.5,\,2.0] để cân bằng Precision--Recall
  \item Chia val theo giống (cultivar-level split) để loại bỏ data leakage
  \item Data augmentation với ảnh không có plastic boards để cải thiện generalization
  \item Nghiên cứu xAI tương thích với query-based Mask2Former (DINO, Transformer rollout)
\end{itemize}

% ─── References ────────────────────────────────────────────────
\bibliographystyle{unsrtnat}
\bibliography{references}

\end{document}
